% ============================== 1. ЦЕЛ И МЕТОДОЛОГИЯ ==============================
\chapter{Цел на изследването и методология за постигане на целта}\label{ch:method}

\begin{keybox}[Цел на изследването]
Да се синтезират, обучат и експериментално сравнят девет модела за откриване на прониквания в мрежовия трафик на център за данни: три атаки (DDoS, BruteForce и Bot), всяка с три алгоритъма (Random Forest, SVM и XGBoost). Моделите се оценяват по TP, TN, FP, FN, Accuracy, Precision и Sensitivity (Recall), допълнени с F1 и Specificity. Методите за обработка и редукция на данните и рискът от преобучение се анализират количествено.
\end{keybox}

Методологията следва цикъла на ML-базирана IDS: избор на атаки и данни, предварителна обработка с редукция на размерността, синтез и настройка на моделите, тестване върху невиждани данни (оперативна фаза) и статистическа оценка~\cite{buczak2016,arp2022}. Всички случайни процеси използват фиксирано начало на генератора (seed~42), така че експериментът се възпроизвежда изцяло.

% ---------------------------------------------------------------------------
\section{Избор на атаки и набори от данни}\label{sec:attacks}

Трите атаки покриват три различни етапа и цели на нападателя в един център за данни.

\emph{Distributed Denial of Service (DDoS)} е шумна, обемна атака: множество атакуващи машини едновременно изпращат трафик, който изчерпва пропускателната способност, таблицата на връзките или ресурсите на приложението~\cite{mirkovic2004}. В използвания набор DDoS е проведена с инструментите HOIC (наводнение с HTTP заявки) и LOIC-UDP (наводнение с UDP пакети). Двете разновидности образуват един клас „DDoS“, защото практическата IDS трябва да разпознава семейството атаки, а не отделния инструмент.

\emph{Brute Force (BruteForce)}, атака по речник, е опит за неоторизиран достъп: нападателят автоматично пробва хиляди двойки потребител--парола срещу услуга за вход. В набора тя е проведена с инструмента Patator срещу FTP и SSH сървъри~\cite{sharafaldin2018}. Всеки опит е кратка, правилно оформена сесия, но опитите се повтарят с еднаква структура.

\emph{Botnet (Bot)} е етапът след успешното проникване. Заразената машина поддържа връзка с команден сървър, получава команди, изпраща снимки на екрана и откраднати данни. В набора е използван Botnet Ares. Трафикът му е с малък обем и периодичен, затова се различава от DDoS по характер.

Изискването е работа \emph{само с реални и достоверни данни}. Затова са сравнени публичните набори с реално уловен трафик (\cref{tab:datasets}).

\begin{table}[H]
\centering\small
\caption{Сравнение на публичните набори от данни за IDS}\label{tab:datasets}
\begin{tabularx}{\textwidth}{@{}l c L c@{}}
\toprule
\textsf{\textbf{Набор}} & \textsf{\textbf{Година}} & \textsf{\textbf{Среда и бележки}} & \textsf{\textbf{Център за данни}}\\
\midrule
KDD Cup 99 / NSL-KDD & 1999/2009 & военна мрежа от 90-те; остарели атаки, много дубликати~\cite{tavallaee2009} & не\\
CIC-IDS2017 & 2017 & корпоративна мрежа, 5 дни реален трафик, достъп след регистрация & частично\\
\rowcolor{boxbg}CSE-CIC-IDS2018 & 2018 & \textbf{облачна инфраструктура в AWS}: 420 машини, 30 сървъра, 10 дни, 14 вида атаки; публичен достъп~\cite{sharafaldin2018,cse2018web} & \textbf{да}\\
CICIoT2023 & 2023 & 105 реални IoT устройства, 33 атаки~\cite{neto2023}; изследван предварително & не\\
\bottomrule
\end{tabularx}
\end{table}

\paragraph{Предварително изследване върху CICIoT2023.} Първоначално изследването беше проведено върху IoT набора CICIoT2023~\cite{neto2023,cic2023web} с атаките DDoS и Port Scan (Recon-PortScan)~\cite{lyon2009}. DDoS се откриваше с Accuracy над 99,9\%, но Port Scan остана около 84\% дори след регуляризация, подбор на признаци, нови признаци и почистване на шума в обучаващите данни (най-добре около 85\%). Причината е в данните: всеки запис описва прозорец от пакети в рамките на един поток, а отделната проба на скенера изглежда като нормална кратка TCP сесия. Затова изследването е пренесено върху CSE-CIC-IDS2018, който е записан в облачен център за данни и е по-подходящ и по тема.

\paragraph{Избраният набор CSE-CIC-IDS2018.} Наборът е създаден съвместно от Communications Security Establishment (CSE) и Canadian Institute for Cybersecurity (University of New Brunswick)~\cite{sharafaldin2018,cse2018web}. Атакуваната организация е изградена в Amazon Web Services (AWS) с 5 отдела, 420 машини и 30 сървъра, а атакуващата инфраструктура е от 50 машини. Записани са 10 работни дни трафик (14.02--02.03.2018). От записите (\texttt{pcap}) инструментът CICFlowMeter-V3 извлича по 80 признака на всеки двупосочен мрежов поток: продължителност, брой и размер на пакетите в двете посоки, интервали между пакетите, TCP флагове, размер на TCP прозореца, периоди на активност и бездействие. Наборът е публикуван в каталога Registry of Open Data on AWS и се изтегля без регистрация~\cite{aws2018}. Той е сред най-използваните в изследванията на IDS върху големи данни~\cite{leevy2020}.

Използвани са три дневни файла (\cref{tab:class-counts}): 21.02.2018 (DDoS), 14.02.2018 (BruteForce) и 02.03.2018 (Bot), общо \val{rows-read} записа. Нормалният трафик от трите дни се обединява. Етикетът на всеки запис е в колоната \texttt{Label}. Две колони не се използват като признаци. \texttt{Timestamp} е моментът на улавяне, а не свойство на трафика. \texttt{Dst Port} е номерът на порта на жертвата. Той би позволил на модела да „запомни“ порта на атакуваната услуга (21 за FTP, 22 за SSH), вместо да научи поведението на потока. Остават \val{n-features} числови признака. Синтетични, генерирани или аугментирани записи не се използват никъде, включително в автоматичните тестове на кода.

\begin{table}[H]
\centering\small
\caption{Изтеглени записи по етикет преди и след почистване}\label{tab:class-counts}
\fitbox{%
\begin{tabular}{@{}lllrrr@{}}
\toprule
\headrow \hd{Набор} & \hd{Етикет} & \hd{Задача} & \hd{Изтеглени} & \hd{След почистване} & \hd{Премахнати}\\
\midrule
CSE-CIC-IDS2018 & \texttt{Benign} & нормален (DDoS, BruteForce) & 1\,790\,843 & 1\,472\,334 & 17,8\%\\
\rowcolor{rowalt}CSE-CIC-IDS2018 & \texttt{DDOS attack-HOIC} & DDoS & 686\,012 & 198\,861 & 71,0\%\\
CSE-CIC-IDS2018 & \texttt{DDOS attack-LOIC-UDP} & DDoS & 1\,730 & 1\,730 & 0,0\%\\
\rowcolor{rowalt}CSE-CIC-IDS2018 & \texttt{FTP-BruteForce} & BruteForce & 193\,360 & 53 & 100,0\%\\
CSE-CIC-IDS2018 & \texttt{SSH-Bruteforce} & BruteForce & 187\,589 & 94\,048 & 49,9\%\\
\rowcolor{rowalt}CIC-IDS2017 & \texttt{BENIGN} & нормален (PortScan) & 979\,288 & 893\,304 & 8,8\%\\
CIC-IDS2017 & \texttt{Portscan} & PortScan & 159\,066 & 1\,683 & 98,9\%\\
\rowcolor{rowalt}двата набора & други етикети & не се използват & 563\,659 & -- & --\\
\textbf{Общо} &  &  & \textbf{4\,561\,547} & \textbf{2\,662\,013} & \\
\bottomrule
\end{tabular}}%

\end{table}
\interprettab{Всеки ред е един етикет от оригиналните файлове. Най-много записи губят атаките BruteForce: \texttt{FTP-BruteForce} съдържа почти само еднакви потоци (неуспешни опити за вход с еднаква структура), затова след премахването на дубликатите от него остават само няколко десетки уникални записа. Нормалният трафик също съдържа много повтарящи се кратки потоци. Дубликатите се премахват, за да не попадне един и същ запис едновременно в обучаващите и в тестовите данни.}

Признаците образуват седем семейства: обем (брой и размер на пакетите в двете посоки, подпотоци), дължини на пакетите (минимум, максимум, средна, стандартно отклонение), интензивност (байтове и пакети в секунда), време (продължителност на потока, интервали между пакетите IAT, периоди на активност и бездействие), TCP флагове (броячи на FIN, SYN, RST, PSH, ACK, URG, ECE), заглавка (протокол, дължина на заглавките, начален размер на TCP прозореца) и насрещни блокове (bulk)~\cite{sharafaldin2018}. Пълният речник на признаците е в хранилището на проекта (\texttt{results/feature\_dictionary.md})~\cite{repo}.

% ---------------------------------------------------------------------------
\section{Предварителна обработка на данните}\label{sec:preprocessing}

Нека суровите данни са $\mathcal{D}=\{(\vect{x}_i, \ell_i)\}_{i=1}^{N}$, където $\vect{x}_i\in\R^{d}$ е векторът на признаците ($d=\val{n-features}$), а $\ell_i$ е оригиналният етикет. Обработката е последователност от трансформации, които превръщат $\mathcal{D}$ в обучаващи и тестови данни за всяка бинарна задача~\cite{arp2022}.

\paragraph{Почистване (премахване на шума).} Премахват се записите с липсващи или безкрайни стойности и пълните дубликати:
\begin{equation}\label{eq:clean}
\mathcal{D}' = \operatorname{unique}\bigl\{(\vect{x}_i,\ell_i)\in\mathcal{D} \;:\; x_{ij}\in\R \ \ \forall j=1,\dots,d\bigr\}.
\end{equation}
Дубликатите се премахват \emph{преди} разделянето на обучаващи и тестови данни, иначе един и същ запис може да попадне в двете части и да надцени качеството~\cite{engelen2021,liu2022}.

\paragraph{Бинарни задачи и почистване на ниво задача.} За атака $a$ с множество от оригинални етикети $\mathcal{L}_a$ (напр. \texttt{FTP-BruteForce} и \texttt{SSH-Bruteforce} за BruteForce) бинарният етикет е
\begin{equation}\label{eq:binary}
y_i = \mathbf{1}\left[\ell_i\in\mathcal{L}_a\right],\qquad i\in\mathcal{I}_a=\{i:\ \ell_i\in\mathcal{L}_a\ \lor\ \ell_i=\txt{Benign}\}.
\end{equation}
Атаката е положителният клас ($y=1$), а нормалният трафик е отрицателният ($y=0$). Някои вектори на признаците се срещат и с етикет на атака, и с етикет на нормален трафик. Това е неразрешимо противоречие (шум в етикетите), известно за наборите на CIC~\cite{liu2022}. Затова в рамките на всяка задача се запазва по едно копие от всеки вектор, а всички вектори с противоречиви етикети се премахват:
\begin{equation}\label{eq:conflict}
\mathcal{C}_a=\bigl\{\vect{x}:\ \exists\,i,k\in\mathcal{I}_a,\ \vect{x}_i=\vect{x}_k=\vect{x},\ y_i\neq y_k\bigr\},\qquad \mathcal{I}_a'=\operatorname{unique}\bigl\{i\in\mathcal{I}_a:\ \vect{x}_i\notin\mathcal{C}_a\bigr\}.
\end{equation}
Така нито един вектор не може да се появи едновременно в обучаващите и в тестовите данни.

\paragraph{Балансиране и разделяне.} От всеки клас се избират случайно без връщане $n_a=\min(|\mathcal{I}_a^{+}|,|\mathcal{I}_a^{-}|,\,10^5)$ записа, така че задачата е балансирана 1:1. Това предпазва моделите от тривиалното решение „всичко е нормално“ и прави Accuracy пригодна метрика, защото при симетрични данни FP и FN имат сравнима тежест~\cite{borovska2022}. Реалистичното небалансирано съотношение се изследва отделно (експеримент E2). Извадката се разделя стратифицирано:
\begin{equation}\label{eq:split}
\mathcal{S}_a=\mathcal{T}_a\,\dot\cup\,\mathcal{V}_a,\qquad |\mathcal{T}_a| = 0{,}7\,|\mathcal{S}_a|,\qquad \tfrac{1}{|\mathcal{T}_a|}\textstyle\sum_{i\in\mathcal{T}_a}y_i=\tfrac{1}{|\mathcal{V}_a|}\sum_{i\in\mathcal{V}_a}y_i=\tfrac12,
\end{equation}
където $\mathcal{T}_a$ е обучаващото множество (фаза на машинно обучение), а $\mathcal{V}_a$ е тестовото (оперативна фаза).

\paragraph{Редукция, етап 1: филтър по дисперсия.} Признаците с дисперсия $s_j^2\le 10^{-8}$ върху обучаващите данни се премахват като константни, където
\begin{equation}\label{eq:var}
s_j^2=\frac{1}{n-1}\sum_{i=1}^{n}\left(x_{ij}-\bar{x}_j\right)^2,\qquad \bar{x}_j=\frac{1}{n}\sum_{i=1}^{n}x_{ij}.
\end{equation}
В CSE-CIC-IDS2018 това са \val{lowvar-DDoS} признака (броячите на насрещни блокове и някои флагове), които CICFlowMeter-V3 записва винаги с нула. Остават $d_f=\val{n-filtered}$ признака.

\paragraph{Редукция, етап 2: подбор на признаци по важност.} Вторият етап избира $k=\val{n-selected}$ оригинални признака, които носят най-много информация за задачата. Върху обучаващите данни се обучава помощен ансамбъл от градиентно усилени дървета (\cref{eq:xgb-obj}), а важността на признак $j$ е средното подобрение на целевата функция (gain) при разклоненията по него:
\begin{equation}\label{eq:gain-imp}
I_j=\frac{1}{|\mathcal{R}_j|}\sum_{r\in\mathcal{R}_j}\operatorname{Gain}(r),\qquad \tilde I_j=\frac{I_j}{\sum_{m}I_m},\qquad \mathcal{J}^{*}=\operatorname{top}_k\bigl\{\tilde I_j\bigr\},
\end{equation}
където $\mathcal{R}_j$ е множеството от разклонения по признак $j$, а $\operatorname{Gain}$ е дефиниран в \cref{eq:xgb-gain}. Подборът запазва оригиналните, интерпретируеми признаци. Той отчита връзката с етикета, докато PCA търси само направленията с голяма дисперсия, които не са непременно полезни за класификацията~\cite{guyon2003}. В предварителното изследване подборът подобри точността при Port Scan с около 3 пункта спрямо PCA. Затова той е основният метод, а PCA е контролен вариант (експеримент E5).

\paragraph{Стандартизация.} Признаците имат много различни мащаби (микросекунди, байтове, броячи), а SVM и PCA са чувствителни към мащаба. Затова всеки признак се стандартизира с параметрите от обучаващите данни:
\begin{equation}\label{eq:zscore}
z_{ij}=\frac{x_{ij}-\bar{x}_j}{s_j}.
\end{equation}

\paragraph{Контролен метод: анализ на главните компоненти (PCA).} PCA търси ортогонални направления с максимална дисперсия~\cite{jolliffe2016}. За стандартизираната матрица $\vect{Z}\in\R^{n\times d_f}$ ковариационната матрица и нейното спектрално разлагане са
\begin{equation}\label{eq:cov}
\bm{\Sigma}=\frac{1}{n-1}\vect{Z}^{\top}\vect{Z},\qquad \bm{\Sigma}\,\vect{w}_k=\lambda_k\vect{w}_k,\qquad \lambda_1\ge\lambda_2\ge\dots\ge\lambda_{d_f}\ge 0.
\end{equation}
Запазва се най-малкият брой компоненти, който обяснява поне 95\% от дисперсията, и всеки запис се проектира върху тях:
\begin{equation}\label{eq:kstar}
\eta(k)=\frac{\sum_{m=1}^{k}\lambda_m}{\sum_{m=1}^{d_f}\lambda_m},\qquad k^{*}=\min\{k:\ \eta(k)\ge 0{,}95\},\qquad \tilde{\vect{x}}=\left[\vect{w}_1,\dots,\vect{w}_{k^*}\right]^{\top}\vect{z}.
\end{equation}

\begin{notebox}[Предпазване от изтичане на информация]
Филтърът, подборът на признаци и стандартизацията са стъпки на един \texttt{sklearn.Pipeline} заедно с класификатора~\cite{pedregosa2011}. Параметрите им, включително важността на признаците, се оценяват \emph{само} върху обучаващите данни, и то вътре във всеки фолд на кръстосаната валидация. Тестовите данни се трансформират с вече научените параметри.
\end{notebox}

% ---------------------------------------------------------------------------
\section{Избор и дизайн на модели за машинно обучение}\label{sec:models}

Избрани са три алгоритъма от различни семейства. Random Forest усреднява много независими дървета (bagging). XGBoost строи дърветата последователно, като всяко поправя грешките на предишните (boosting). SVM с RBF ядро е метод на максималния марж, който строи гладка нелинейна граница. Различните индуктивни пристрастия правят сравнението информативно~\cite{hastie2009}.

\paragraph{Решаващо дърво и Random Forest.} Решаващото дърво (CART) разделя пространството рекурсивно с прагови правила $x_j\le t$. Във всеки възел се избира разделянето, което най-много намалява индекса на Джини~\cite{breiman1984}:
\begin{equation}\label{eq:gini}
G(S_v)=1-\sum_{c\in\{0,1\}}p_{v,c}^2,\qquad p_{v,c}=\frac{|\{i\in S_v:\ y_i=c\}|}{|S_v|}.
\end{equation}
Единичното дърво има висока дисперсия. Random Forest~\cite{breiman2001} обучава $B$ дървета, всяко върху bootstrap извадка от обучаващите данни, като при всяко разделяне разглежда само случайно подмножество от $m$ признака. Вероятността за атака и решението на ансамбъла са
\begin{equation}\label{eq:rf}
\hat{p}(\vect{x})=\frac{1}{B}\sum_{b=1}^{B}\hat{p}_b(\vect{x}),\qquad \hat{y}(\vect{x})=\mathbf{1}\left[\hat{p}(\vect{x})>0{,}5\right],
\end{equation}
където $\hat{p}_b(\vect{x})$ е делът на атаките в листото на дърво $b$. Ако дърветата имат дисперсия $\sigma^2$ и средна двойкова корелация $\rho$, дисперсията на средното им е~\cite{hastie2009}
\begin{equation}\label{eq:rf-var}
\operatorname{Var}\bigl[\hat{p}\bigr]=\rho\,\sigma^2+\frac{1-\rho}{B}\,\sigma^2,
\end{equation}
т.е. ансамбълът е толкова по-стабилен, колкото по-слабо корелирани са дърветата. Минималният брой записи в листо и максималната дълбочина ограничават сложността на всяко дърво и така намаляват преобучението.

\paragraph{Градиентно усилване (XGBoost).} XGBoost~\cite{chen2016} е адитивен модел от $K$ регресионни дървета $f_k$, построени последователно~\cite{friedman2001}. Оценката (логит) и вероятността за атака са
\begin{equation}\label{eq:xgb-model}
F(\vect{x})=\sum_{k=1}^{K}\nu\,f_k(\vect{x}),\qquad \hat{p}(\vect{x})=\sigma\bigl(F(\vect{x})\bigr)=\frac{1}{1+e^{-F(\vect{x})}},
\end{equation}
където $\nu$ е скоростта на обучение. Всяко ново дърво минимизира регуляризираната целева функция
\begin{equation}\label{eq:xgb-obj}
\mathcal{L}^{(k)}=\sum_{i=1}^{n}l\bigl(y_i,\ F^{(k-1)}(\vect{x}_i)+f_k(\vect{x}_i)\bigr)+\Omega(f_k),\qquad \Omega(f)=\gamma T+\tfrac12\lambda\sum_{t=1}^{T}w_t^2,
\end{equation}
където $l$ е логистичната загуба, $T$ е броят листа, а $w_t$ са стойностите в листата. Членът $\Omega$ наказва сложните дървета и е вградената защита срещу преобучение. С приближение от втори ред с градиентите $g_i=\partial l/\partial F$ и вторите производни $h_i=\partial^2 l/\partial F^2$ оптималната стойност на листо и подобрението от разклонение на множеството $S$ на $S_L$ и $S_R$ са
\begin{equation}\label{eq:xgb-gain}
w_t^{*}=-\frac{\sum_{i\in S_t}g_i}{\sum_{i\in S_t}h_i+\lambda},\qquad
\operatorname{Gain}=\frac12\left[\frac{G_L^2}{H_L+\lambda}+\frac{G_R^2}{H_R+\lambda}-\frac{(G_L+G_R)^2}{H_L+H_R+\lambda}\right]-\gamma,
\end{equation}
където $G$ и $H$ са сумите на $g_i$ и $h_i$ в съответния възел. Допълнително всяко дърво се обучава върху случайни 80\% от записите и 80\% от признаците (subsample, colsample), а параметърът \texttt{min\_child\_weight} забранява листа с малка сума на $h_i$.

\paragraph{Метод на опорните вектори (SVM).} SVM търси хиперравнината $\vect{w}^{\top}\phi(\vect{x})+b=0$, която разделя класовете с максимален марж~\cite{cortes1995}. С етикети $t_i=2y_i-1\in\{-1,+1\}$ задачата с мек марж е
\begin{equation}\label{eq:svm-primal}
\min_{\vect{w},b,\bm{\xi}}\ \frac12\|\vect{w}\|^2+C\sum_{i=1}^{n}\xi_i\qquad\txt{при}\qquad t_i\left(\vect{w}^{\top}\phi(\vect{x}_i)+b\right)\ge 1-\xi_i,\quad \xi_i\ge 0.
\end{equation}
Параметърът $C$ регулира компромиса между широк марж и малко грешки. Дуалната задача съдържа данните само чрез скаларни произведения, които се заменят с ядро~\cite{scholkopf2002}:
\begin{equation}\label{eq:svm-dual}
\max_{\bm{\alpha}}\ \sum_{i=1}^{n}\alpha_i-\frac12\sum_{i=1}^{n}\sum_{k=1}^{n}\alpha_i\alpha_k t_i t_k K(\vect{x}_i,\vect{x}_k)\qquad\txt{при}\qquad 0\le\alpha_i\le C,\ \ \sum_{i}\alpha_i t_i=0.
\end{equation}
Използва се радиалното базисно ядро, а решаващата функция зависи само от опорните вектори ($\alpha_i>0$):
\begin{equation}\label{eq:rbf}
K(\vect{x},\vect{x}')=\exp\left(-\gamma\,\|\vect{x}-\vect{x}'\|^2\right),\qquad f(\vect{x})=\sum_{i:\ \alpha_i>0}\alpha_i t_i K(\vect{x}_i,\vect{x})+b.
\end{equation}
Малко $\gamma$ дава гладка граница, а голямо $\gamma$ дава локална граница с риск от преобучение. Обучението (SMO в LIBSVM~\cite{chang2011}) има сложност между $O(n^2)$ и $O(n^3)$, а класификацията на един запис е $O(n_{SV}\cdot d)$. Затова SVM се обучава върху ограничена извадка до \num{50000} записа.

\paragraph{Избор на хиперпараметри и защита от преобучение.} Хиперпараметрите $\bm{\theta}$ се избират с обхождане на решетка и $K=5$-кратна стратифицирана кръстосана валидация~\cite{kohavi1995}. За всяка конфигурация се изчисляват средната оценка и стандартната грешка по фолдовете:
\begin{equation}\label{eq:cv}
\overline{\operatorname{F1}}(\bm{\theta})=\frac{1}{K}\sum_{k=1}^{K}\operatorname{F1}\Bigl(h_{\bm{\theta}}^{(\mathcal{T}\setminus\mathcal{F}_k)},\ \mathcal{F}_k\Bigr),\qquad \operatorname{SE}(\bm{\theta})=\frac{s_{\operatorname{F1}}(\bm{\theta})}{\sqrt{K}},
\end{equation}
където $h_{\bm{\theta}}^{(\mathcal{A})}$ е целият конвейер, обучен върху $\mathcal{A}$, а $\mathcal{F}_k$ е $k$-тият фолд. Вместо конфигурацията с най-висока оценка се избира \emph{най-простата} конфигурация, чиято оценка е в рамките на една стандартна грешка от най-добрата (правилото „една стандартна грешка“~\cite{breiman1984,hastie2009}):
\begin{equation}\label{eq:one-se}
\bm{\theta}^{*}=\argmin_{\bm{\theta}\in\Theta}\Bigl\{c(\bm{\theta}):\ \overline{\operatorname{F1}}(\bm{\theta})\ \ge\ \overline{\operatorname{F1}}(\hat{\bm{\theta}})-\operatorname{SE}(\hat{\bm{\theta}})\Bigr\},\qquad \hat{\bm{\theta}}=\argmax_{\bm{\theta}\in\Theta}\overline{\operatorname{F1}}(\bm{\theta}),
\end{equation}
където $c(\bm{\theta})$ е мярка за сложност: по-голям минимален брой записи в листо и по-малка дълбочина за дърветата, по-малки $C$ и $\gamma$ за SVM. Конфигурациите в рамките на една стандартна грешка са статистически неразличими, а по-простата обобщава по-добре. Критерият F1 балансира пропуснатите атаки и фалшивите аларми (\cref{eq:f1}). Преобучението се проверява явно: за всеки модел се сравняват точността върху обучаващите данни, оценката от кръстосаната валидация и точността върху тестовите данни.

% ---------------------------------------------------------------------------
\section{Експериментална оценка на моделите}\label{sec:metrics}

Ефективността на моделите се оценява чрез стандартните метрики за качество на класификацията в системите за откриване на проникване: TP, TN, FP, FN, Precision, Accuracy и Sensitivity (Recall). Те показват както общата точност на модела, така и доколко той открива реалните атаки и избягва фалшивите аларми. Резултатът от класификацията на тестовите записи се обобщава в матрица на объркване (\cref{tab:cm}). Положителният клас е атаката. По аналогия с медицинската диагноза~\cite{borovska2022} TP и TN са желаните решения, FP е фалшива тревога, а FN (пропусната атака) е най-нежеланата грешка, защото атаката продължава незабелязано.

\begin{table}[H]
\centering\small
\caption{Матрица на объркване: определение на TP, TN, FP и FN}\label{tab:cm}
\begin{tabular}{@{}l l c c@{}}
\toprule
 & & \multicolumn{2}{c}{\textsf{\textbf{Предсказан клас}}}\\
\cmidrule(l){3-4}
 & & Атака ($\hat y=1$) & Нормален ($\hat y=0$)\\
\midrule
\multirow{2}{*}{\textsf{\textbf{Реален клас}}} & Атака ($y=1$) & \cellcolor{boxbg}TP: открита атака & FN: пропусната атака\\
 & Нормален ($y=0$) & FP: фалшива аларма & \cellcolor{boxbg}TN: правилно пропуснат\\
\bottomrule
\end{tabular}
\end{table}

Метриките са изброени в реда на учебния материал~\cite{borovska2022}. Всяка отговаря на различен въпрос и се тълкува „на 100 случая“.

\paragraph{Акуратност (Accuracy)} е делът на всички верни решения:
\begin{equation}\label{eq:accuracy}
\operatorname{Accuracy}=\frac{TP+TN}{TP+TN+FP+FN}.
\end{equation}
Accuracy~=~0,99 означава, че от 100 тестови записа 99 са класифицирани правилно. Метриката е най-интуитивна, но е подходяща главно за симетрични данни, при които FP и FN са близки.

\paragraph{Прецизност (Precision)} е делът на истинските атаки сред всички аларми:
\begin{equation}\label{eq:precision}
\operatorname{Precision}=\frac{TP}{TP+FP}.
\end{equation}
Precision~=~0,95 означава, че от 100 аларми 95 са истински атаки, а 5 са фалшиви. В примера на материала за Precision е описано тълкуване на Recall. Тук е използвано коректното тълкуване: делът на верните сред положителните решения~\cite{fawcett2006}.

\paragraph{Чувствителност (Sensitivity, Recall)} е делът на откритите атаки:
\begin{equation}\label{eq:sensitivity}
\operatorname{Sensitivity}=\operatorname{Recall}=\operatorname{TPR}=\frac{TP}{TP+FN}.
\end{equation}
Recall~=~0,97 означава, че от 100 атаки 97 се откриват, а 3 се пропускат. Тъй като FN е най-нежеланата грешка, Recall е водещата метрика за сигурността.

\paragraph{F1 оценка} е хармоничното средно на Precision и Recall:
\begin{equation}\label{eq:f1}
\operatorname{F1}=2\cdot\frac{\operatorname{Precision}\cdot\operatorname{Recall}}{\operatorname{Precision}+\operatorname{Recall}}=\frac{2\,TP}{2\,TP+FP+FN}.
\end{equation}
F1 е висока само когато и двете съставки са високи. Затова тя е предпочитаната обобщена метрика при несиметрични данни, при които Accuracy подвежда~\cite{borovska2022,saito2015}.

\paragraph{Специфичност (Specificity)} е делът на нормалния трафик, който не предизвиква аларма:
\begin{equation}\label{eq:spec}
\operatorname{Specificity}=\operatorname{TNR}=\frac{TN}{TN+FP},\qquad \operatorname{FPR}=1-\operatorname{Specificity}.
\end{equation}
Specificity~=~0,98 означава, че от 100 нормални записа 98 минават без аларма, а 2 предизвикват фалшива тревога.

\paragraph{ROC крива и реална честота на атаките.} Трите модела дават непрекъсната оценка $s(\vect{x})$ ($\hat p$ при RF и XGBoost и $f$ при SVM). При праг $\theta$ решението е $\hat y=\mathbf{1}[s(\vect{x})\ge\theta]$. ROC кривата съдържа точките $(\operatorname{FPR}(\theta),\operatorname{TPR}(\theta))$ за всички прагове, а площта под нея е вероятността случайна атака да получи по-висока оценка от случаен нормален запис~\cite{fawcett2006}:
\begin{equation}\label{eq:auc}
\operatorname{AUC}=\int_0^1\operatorname{TPR}\,d\operatorname{FPR}=P\bigl(s(\vect{x}^{+})>s(\vect{x}^{-})\bigr).
\end{equation}
Precision зависи и от дела $\pi$ на атаките в трафика. От теоремата на Бейс следва~\cite{axelsson2000}
\begin{equation}\label{eq:baserate}
\operatorname{Precision}(\pi)=\frac{\operatorname{TPR}\cdot\pi}{\operatorname{TPR}\cdot\pi+\operatorname{FPR}\cdot(1-\pi)}.
\end{equation}
При малко $\pi$ дори нисък FPR дава повече фалшиви, отколкото истински аларми. Това е \emph{проблемът на базовата честота}, изследван в експеримент E2 с Precision--Recall криви~\cite{saito2015}. Там се избират и работни прагове: най-високият праг, при който Recall достига целева стойност (0,90 или 0,95), за да се ограничат пропуснатите атаки.

% ---------------------------------------------------------------------------
\section{Графично представяне и анализ на резултатите}\label{sec:stats}

Резултатите се анализират в два слоя: статистическа проверка на надеждността на метриките и графично представяне на всички етапи на експеримента~\cite{efron1993}.

\paragraph{Bootstrap доверителни интервали.} Несигурността на метриките се оценява с непараметричен bootstrap~\cite{efron1993}: от $n$-те тестови записа се теглят с връщане $R=\num{1000}$ извадки с размер $n$, върху всяка се изчислява метриката $\hat M^{*}_r$, а 95\% перцентилният интервал е
\begin{equation}\label{eq:bootstrap}
\operatorname{CI}_{95\%}(M)=\left[\hat{M}^{*}_{(0{,}025)},\ \hat{M}^{*}_{(0{,}975)}\right].
\end{equation}

\paragraph{Тест на McNemar.} Всяка двойка класификатори A и B се сравнява върху едни и същи тестови записи по записите, по които се различават~\cite{mcnemar1947,dietterich1998}: $b$ е броят записи, верни при A и грешни при B, а $c$ е обратното. При нулевата хипотеза за еднакъв процент грешки
\begin{equation}\label{eq:mcnemar}
\chi^2=\frac{\left(|b-c|-1\right)^2}{b+c}\ \sim\ \chi^2_1,
\end{equation}
а при $b+c<25$ се използва точният биномен тест.

\paragraph{Едномерна разделимост.} За всеки признак се сравняват емпиричните функции на разпределение при атака ($F_1$) и нормален трафик ($F_0$) със статистиката на Колмогоров--Смирнов
\begin{equation}\label{eq:ks}
D_j=\sup_{t\in\R}\ \bigl|F_{1,j}(t)-F_{0,j}(t)\bigr|\in[0,1].
\end{equation}
$D_j=1$ означава, че един праг по признак $j$ разделя класовете идеално~\cite{hastie2009}.

\paragraph{Графично представяне.} Всички графики се генерират автоматично след всяко изпълнение с еднакво цветово кодиране: Random Forest в синьо, SVM в оранжево, XGBoost в зелено, атаката в червено, нормалният трафик в сиво. Те покриват данните и обработката, качеството на моделите, поведението им при обучение и работата при реално съотношение. След всяка фигура в \cref{ch:results} следва обяснение какво е показано~\cite{repo}.
